\uuid{KlB0}
\exo7id{1687}
\titre{exo7 1687}
\auteur{bodin}
\organisation{exo7}
\datecreate{1999-11-01}
\isIndication{false}
\isCorrection{false}
\chapitre{Réduction d'endomorphisme, polynôme annulateur}
\sousChapitre{Trigonalisation}
\cpge{mpsi-mp2i.calcul-matriciel-systemes.e, pcsi.determinants.c, ptsi.determinants.a, mpsi-mp2i.changements-bases.c, bcpst1.matrices, mp-mpi.reduction-endomorphismes.b, pc.reduction.a, psi.reduction-endomorphismes.a, pt.reduction-endomorphismes.a, mp-mpi.reduction-endomorphismes.c, pc.reduction.b, psi.reduction-endomorphismes.b, pt.reduction-endomorphismes.b, mp-mpi.reduction-endomorphismes.e, pc.reduction.e, psi.reduction-endomorphismes.e, pt.reduction-endomorphismes.d, bcpst2.valeurs-propres-vecteurs-propres.a, bcpst2.valeurs-propres-vecteurs-propres.b}
\module{Algèbre}
\niveau{L2}
\difficulte{}

\contenu{
\question{
Mettre sous forme triangulaire les matrices suivantes :
\[  \begin{array}{ll}
\left( \begin{array}{ccc} 4 & 2 & -2 \\ 1 & 5 & -1
\\ 1 & 3 & 1 \end{array}\right); \quad & \qquad
\frac{1}{2} \left( \begin{array}{ccc} 0 & 2 & 2 \\
1 & 3 & -1 \\ -1 & 3 & 3
\end{array}\right). \end{array} \]
}
}
